% !TeX root = ../../main.tex
\subsection{Uniform bounds at the two kernel endpoints}
\label{subsec:b:endpoints}

Here the remainder exponent is a fixed real number
$1/2<\nu<1$. Angular domains are always used in nested pairs:
estimates on a smaller cone are obtained from holomorphic functions
on a larger one. This convention permits Cauchy estimates of any
fixed order, and avoids a loss of derivatives in a finite-layer
induction. The bounds below concern general input functions;
Proposition~\ref{prp:b:common-domains} applies them to the actual
Borel series.

Put $S=R_\theta-x$. We say that a holomorphic input
$\mathcal H(\theta;x,y)$ has \emph{scale-$a$ endpoint bounds},
where $a\geq0$ is real, if it is bounded by $C|S|^a$ on the small-$|S|$
part of $\mathcal C_\gamma$, and has the expansions
\begin{equation}\label{eq:b:endpoint-class}
 \mathcal H
 =H_0^\pm(\theta,S)+\sqrt{S\mp y}\,H_1^\pm(\theta,S)+E^\pm
                 \quad\text{near }y=\pm S,
\end{equation}
with
\begin{equation}\label{eq:b:endpoint-norms}
 |H_0^\pm|\leq C|S|^a,\qquad
 |H_1^\pm|\leq C|S|^{a-1/2},\qquad
 |E^\pm|\leq C|S|^a|1\mp y/S|^\nu.
\end{equation}
The coefficient functions are holomorphic in the slit $S$ sector
and in a complex neighborhood of the compact $\theta$ interval
under consideration. Their bounds hold on slightly larger sectors
and parameter neighborhoods as well. Constants are uniform there.
All statements about this class allow a fixed reduction of those
neighborhoods.

In particular, Cauchy's formula, applied with radius a fixed
multiple of $|S\mp y|$, gives for every fixed integer $k\geq0$
\begin{equation}\label{eq:b:remainder-derivatives}
 |S^k\partial_y^k E^\pm|
       \leq C_k|S|^a|1\mp y/S|^{\nu-k}.
\end{equation}
The coefficient functions have the corresponding bounds after any
fixed number of $S\partial_S$ and $\theta$ derivatives. These are
consequences of the larger-domain value bounds, not extra regularity
assumptions at the endpoints. For example, for every fixed $k\geq1$,
\begin{equation}\label{eq:b:input-derivatives}
 |\partial_x^k\mathcal H|+|\partial_y^k\mathcal H|
 \leq C_k|S|^{a-k}
 \left(1+|1-y/S|^{1/2-k}+|1+y/S|^{1/2-k}\right).
\end{equation}
For $x$ derivatives, differentiate the coefficients and
$S=R_\theta-x$ in \eqref{eq:b:endpoint-class}; the remainder is
estimated on the same nested cones. Mixed derivatives obey the
same bound with their total order in place of $k$.

\begin{lmm}[The linear endpoint integral]\label{lmm:b:linear-endpoint}
Suppose $\mathcal H$ has scale-$a$ endpoint bounds, $a\geq0$.
Then $\mathcal K\mathcal H$ has such bounds with the same scale,
and near the positive endpoint its square-root coefficient is
\begin{equation}\label{eq:b:linear-endpoint}
 \begin{aligned}
 (\mathcal K\mathcal H)(\theta;x,y)
 =K_0^+(\theta,S)
   +\sqrt{S-y}\int_0^1H_1^+(\theta,\tau S)\,\dd\tau+E_K^+,\\
 |E_K^+|\leq C|S|^a|1-y/S|^\nu.
 \end{aligned}
\end{equation}
Here $|K_0^+|\leq C|S|^a$. At the negative endpoint the analogous
square-root coefficient is
$-\int_0^1H_1^-(\theta,\tau S)\,\dd\tau$.
Every part except the $m=1$ branch with the corresponding sign
has a bounded first normalized $y$ derivative at that endpoint.
The remainder estimates hold with any fixed number of derivatives
on smaller angular domains.
\end{lmm}
\begin{prf}
Away from the two endpoint neighborhoods, cone comparison and the
cubic Taylor remainder give the scale bound directly. We concentrate
on $y=S$.
For the positive $m=1$ branch, write
\[
 S'=S-(1-\tau)y,\qquad
 S'-\tau y=S-y.
\]
Cone addition gives
$|S'|\asymp |S|(\tau+|1-y/S|)$ near $y=S$.
Insert \eqref{eq:b:endpoint-class} into this branch. The integral
of its remainder is bounded by
\[
 C|S|^a|1-y/S|^\nu
       \int_0^1(\tau+|1-y/S|)^{a-\nu}\,\dd\tau,
\]
which is uniform because $a-\nu>-1$.

To freeze the square-root coefficient at $y=S$, split the
$\tau$ interval at $|1-y/S|$. On the first part, both the input
and the proposed frozen term have integral
$O(|S|^a|1-y/S|^{a+1})$. On the second part, the coefficient
change is bounded by its $S$ derivative. Its integral is at most
\[
 C|S|^a|1-y/S|^{3/2}
      \int_{|1-y/S|}^1\tau^{a-3/2}\,\dd\tau.
\]
This is $O(|S|^a|1-y/S|^\nu)$. Freezing the constant coefficient
similarly costs at most
$C|S|^a|1-y/S|(1+|\log|1-y/S||)$, which obeys the same bound.
The coefficient integral in \eqref{eq:b:linear-endpoint} converges
at zero because $a-1/2>-1$.

For $m=1$ the other branch can meet this endpoint only at
$\tau=1$. Its first normalized derivative has an integrable bound
$C|S|^a(|1-y/S|+1-\tau)^{-1/2}$.
The subtracted first-order term has the extra factor $1-\tau$;
its derivative is bounded by the integral of
\[
 C|S|^a(1-\tau)
                  (|1-y/S|+1-\tau)^{-3/2},
\]
which is also uniform.

For $m\geq2$ keep the entire cubic
difference together; its individual linear terms must not be
summed separately. Taylor's integral remainder and
\eqref{eq:b:input-derivatives} bound its first normalized derivative,
before the outer factor $m$, by a constant times
\[
 |S|^a m^{-3}(1-\tau)^3
                  (|1-y/S|+1-\tau)^{-7/2}.
\]
The $\tau$ integral is uniform, and the remaining $m^{-2}$ sum
converges. Hence these nonprincipal parts have a finite endpoint
limit and an $O(|S|^a|1-y/S|)$ variation. This proves the value
estimate in \eqref{eq:b:linear-endpoint}.

The negative endpoint
follows by exchanging the two branches, including their opposite
signs. Cauchy's formula on smaller cones then gives every fixed
derivative bound, in particular the third derivative required by
the next transport step.
\end{prf}

\begin{lmm}[The quadratic endpoint integral]\label{lmm:b:quadratic-endpoint}
Let $\mathcal H$ and $\mathcal J$ have endpoint scales $a,c\geq0$.
The function
\begin{equation}\label{eq:b:quadratic-input}
 \mathcal Q=y^2\int_0^1\int_0^1
 t\,\partial_x\mathcal H(\theta;x,t\tau y)
       \partial_x\mathcal J(\theta;x,t(1-\tau)y)
                              \,\dd\tau\,\dd t
\end{equation}
has scale $a+c$. At either endpoint it has a constant term and
a remainder bounded by $C|S|^{a+c}|1\mp y/S|^\nu$, with zero
square-root coefficient. Its first normalized $y$ derivative is
bounded by $C|S|^{a+c}$; its second is bounded by
$C|S|^{a+c}(1+|1\mp y/S|^{-1/2})$.
The value remainder again gives every fixed derivative estimate
on smaller cones.
\end{lmm}
\begin{prf}
The two $x$ derivatives cost $|S|^{-2}$, and the prefactor $y^2$
restores that scale. Cone addition gives
\[
 |S-t\tau y|\geq c_\gamma
       \bigl((1-t\tau)|S|+t\tau|S-y|\bigr),
\]
and the corresponding inequality for all other signs and factors.
Since $t\tau+t(1-\tau)=t\leq1$, the two factors cannot both
approach the same endpoint. In fact at least one of the two
coefficients is at most $1/2$, and its distance is bounded below
by a fixed multiple of $|S|$. The same argument excludes
simultaneous approach to opposite endpoints.

A singular factor can occur only at the corner $(t,\tau)=(1,1)$
or $(1,0)$. For its first normalized $y$ derivative the largest
local bound, apart from $C|S|^{a+c}$, is
\[
 (|1-y/S|+(1-t)+(1-\tau))^{-3/2}
\]
at the first corner, and the corresponding bound at the second.

This is integrable in two variables, uniformly up to the endpoint.
For a second derivative the exponent is $-5/2$, whose integral is
$O(1+|1-y/S|^{-1/2})$. The remainder terms with exponent $\nu$
are less singular because $\nu>1/2$.

Thus the first derivative
is bounded and the endpoint value exists, with a remainder
$O(|S|^{a+c}|1-y/S|)$. It is therefore also $O(|S|^{a+c}|1-y/S|^\nu)$.
Normal convergence gives holomorphy; larger-domain Cauchy estimates
give the remaining assertions.
\end{prf}
